% PAGES: 298-302
% Section 9.8 "Exercises" (folios 282-286), the final section of chapter 9.
% This is the LAST file of the LAST chunk of chapter 9: per the standing
% brief, the chapter's whole exercise list (15 exercises, folios 282-286)
% goes in this ONE file inside a single \begin{exercises}...\end{exercises}
% environment, since that environment opens a fresh enumerate and splitting
% it across files would restart the numbering at 1.
%
% Chapter 10 ("Mathematical appendix and technical results") begins cleanly
% on the very next PDF page (303, folio 287, rendered and confirmed via its
% own CropBox) immediately after exercise 15 ends, so nothing is left open
% at the end of this file, and nothing is left open at the end of the
% chapter.

\section{Exercises}

\begin{exercises}

\item Two stocks trade at \$100 and \$60, respectively. Their three--months
returns have expected values of 8\% and 4\%, respectively, and standard
deviation of 20\% and 15\%, respectively. The correlation of the returns is
25\%.
\begin{subparts}
\item Consider a portfolio made of 150 shares of the first stock and 500
shares of the second stock. What are the weights of each stock in this
portfolio?
\item Assume that you have \$500,000 to invest. Find a portfolio made of the
two stocks that has a 9\% expected return.
\item Identify the two portfolios fully invested in the two assets that have
a 14.5\% standard deviation of return. What are the expected returns of the
two portfolios?
\end{subparts}

\item Assume that the asset allocation for a \$100 million maximum return
portfolio invested in three assets is \$30 million in the first asset, \$10
million in the second asset, \$40 million in the third asset, and \$20
million in cash. What is the asset allocation of the tangency portfolio made
of these three assets?

\item Consider three assets with the following expected values, standard
deviations, and correlations of their returns:
\[
\begin{aligned}
\mu_1 &= 0.08; & \sigma_1 &= 0.25; & \rho_{1,2} &= -0.25;\\
\mu_2 &= 0.12; & \sigma_2 &= 0.25; & \rho_{2,3} &= -0.25;\\
\mu_3 &= 0.16; & \sigma_3 &= 0.30; & \rho_{1,3} &= 0.25.
\end{aligned}
\]
The risk--free interest rate is 4\%.
\begin{subparts}
\item Find the asset allocation corresponding to the tangency portfolio.
What are the expected value and the standard deviation of the return of the
tangency portfolio?
\item Find the asset allocation corresponding to the minimum variance
portfolio. What are the expected value and the standard deviation of the
return of the minimum variance portfolio?
\end{subparts}

\item Assume that you invest \$10 million in two different assets and cash.
The three--months returns of the two assets have expected values of 8\% and
12\%, respectively, and standard deviations of 15\% and 20\%, respectively.
The correlation of the returns of the two assets is 25\%. The risk--free
interest rate is 5\%.
\begin{subparts}
\item Find the asset allocation for the tangency portfolio.
\item Find the asset allocation for a minimum variance portfolio with 7\%
expected return, and the standard deviation of the return of this portfolio.
\item Find the asset allocation for a minimum variance portfolio with 11\%
expected return, and the standard deviation of the return of this portfolio.
\item Find the asset allocation for a maximum return portfolio with 12\%
standard deviation of return, and the expected return of this portfolio.
\item Find the asset allocation for a maximum return portfolio with 18\%
standard deviation of return, and the expected return of this portfolio.
\item Assume that the risk--free interest rate changes to 5.25\%. How do you
adjust the asset allocation of the minimum variance portfolio with 7\%
expected return in order to maintain a minimum variance portfolio with 7\%
expected return?
\end{subparts}

\item Consider two assets with three--months returns with expected values of
6\% and 10\%, respectively, and standard deviations of 25\% and 35\%,
respectively. The correlation of the three--months returns of the two assets
is $-25\%$. The risk--free interest rate is 3\%. Assume that you invest \$100
million in the two assets.
\begin{subparts}
\item Find the minimum variance of the return of a portfolio made of the two
assets. What is the expected return of this portfolio?
\item What is the 10--day 98\% VaR of a \$100 million portfolio, if it is
invested in the first asset, invested in the second asset, or invested in
the minimum variance portfolio, respectively?
\end{subparts}

\item Consider three assets with the following expected values, standard
deviations, and correlations of their returns:
\[
\begin{aligned}
\mu_1 &= 0.06; & \sigma_1 &= 0.18; & \rho_{1,2} &= -0.50;\\
\mu_2 &= 0.09; & \sigma_2 &= 0.20; & \rho_{2,3} &= -0.25;\\
\mu_3 &= 0.12; & \sigma_3 &= 0.24; & \rho_{1,3} &= 0.15.
\end{aligned}
\]
The risk--free interest rate is 3\%.
\begin{subparts}
\item Find the asset allocation for the tangency portfolio. What are the
expected return and the standard deviation of the return of the tangency
portfolio?
\item Find the asset allocation for a minimum variance portfolio with 10\%
expected return, and the standard deviation of the return of this portfolio.
\item Find the asset allocation for a maximum return portfolio with 20\%
standard deviation of return, and the expected return of this portfolio.
\end{subparts}

\item Consider five assets with the following expected values, standard
deviations, and correlations of their returns:
\[
\begin{aligned}
\mu_1 &= 0.08; & \sigma_1 &= 0.25;\\
\mu_2 &= 0.12; & \sigma_2 &= 0.25;\\
\mu_3 &= 0.16; & \sigma_3 &= 0.30;\\
\mu_4 &= 0.18; & \sigma_4 &= 0.32;\\
\mu_5 &= 0.21; & \sigma_5 &= 0.35.
\end{aligned}
\]
\[
\begin{aligned}
\rho_{1,2} &= -0.25; & \rho_{1,3} &= -0.25; & \rho_{1,4} &= 0.35; & \rho_{1,5} &= -0.10;\\
\rho_{2,3} &= 0.30; & \rho_{2,4} &= -0.50; & \rho_{2,5} &= 0.10;\\
\rho_{3,4} &= -0.30; & \rho_{3,5} &= -0.35; & \rho_{4,5} &= 0.65;
\end{aligned}
\]
The risk--free interest rate is 5\%.
\begin{subparts}
\item Find the asset allocation for a minimum variance portfolio with 17\%
expected return, and the standard deviation of the return of this portfolio;
\item Find the asset allocation for a maximum return portfolio with 30\%
standard deviation of return, and the expected return of this portfolio.
\end{subparts}

\item Consider four assets with the following expected returns over a fixed
time period:
\[
\mu_1 = 4\%;\quad \mu_2 = 3.5\%;\quad \mu_3 = 5\%;\quad \mu_4 = 3.4\%,
\]
and with the following covariance matrix of their returns over the same time
period:
\[
\begin{pmatrix}
0.09 & 0.01 & 0.03 & -0.015\\
0.01 & 0.0625 & -0.02 & -0.01\\
0.03 & -0.02 & 0.1225 & 0.02\\
-0.015 & -0.01 & 0.02 & 0.0576
\end{pmatrix}.
\]
Assume that the risk--free interest rate is 1\%.
\begin{subparts}
\item Find the asset allocation for the tangency portfolio. Find the
expected value and the standard deviation of the return of the tangency
portfolio. What is the Sharpe ratio of the tangency portfolio?
\item Find the asset allocation for a minimum variance portfolio with 3\%
expected return, and the standard deviation of the return of this portfolio.
What is the Sharpe ratio of this portfolio?
\item Find the asset allocation for a maximum return portfolio with 27\%
standard deviation of return, and the expected return of this portfolio.
What is the Sharpe ratio of this portfolio?
\item Find the asset allocation for the minimum variance portfolio fully
invested in the assets (i.e., with no cash position). What is the Sharpe
ratio of this portfolio?
\end{subparts}

\item Consider three assets with the following expected values, standard
deviations, and correlations of their returns:
\[
\begin{aligned}
\mu_1 &= 0.05; & \sigma_1 &= 0.15; & \rho_{1,2} &= -0.25;\\
\mu_2 &= 0.09; & \sigma_2 &= 0.20; & \rho_{2,3} &= 0.25;\\
\mu_3 &= 0.10; & \sigma_3 &= 0.25; & \rho_{1,3} &= 0.50.
\end{aligned}
\]
The risk--free interest rate is 2\%.
\begin{subparts}
\item Find the asset allocation for a minimum variance portfolio with 8\%
expected return, and the standard deviation of the return of this portfolio.
\item Assume that the returns of the three assets have a joint multivariate
normal distribution. Find the probability density function of the return of
the minimum variance portfolio with 8\% expected return.
\item Find the probability that the return of the minimum variance portfolio
with 8\% expected return is between 7\% and 9\%. Also, find the probability
that the return of this portfolio is below 5\%, and the probability that the
return of this portfolio is above 10\%.
\item Consider a portfolio equally invested in each of the three assets.
Note that the expected return of this portfolio is 8\%. Find the
probabilities that the return of this portfolio is between 7\% and 9\%, is
below 5\%, and is above 10\%, respectively.
\end{subparts}

\item Let $\mu_1,\sigma_1$ and $\mu_2,\sigma_2$ be the expected values and
the standard deviations of the returns of two assets over a fixed time
period, respectively, and let $\rho_{1,2}$ be the correlation of the returns
of the two assets over the same time period. Recall that the minimum
variance portfolio fully invested in two assets is obtained by allocating
\[
w_1 \;=\; \frac{\sigma_2(\sigma_2-\rho_{1,2}\sigma_1)}{(\sigma_1^2 - 2\sigma_1\sigma_2\rho_{1,2}+\sigma_2^2)}
\]
of the portfolio value to the first asset and $w_2 = 1-w_1$ of the portfolio
value to the second asset.

Show that the minimum variance portfolio has only long asset positions,
i.e., $0 \le w_1,w_2 \le 1$, if and only if
\[
\rho_{1,2} \;\le\; \min\left(\frac{\sigma_1}{\sigma_2},\frac{\sigma_2}{\sigma_1}\right).
\]

\item (i) Show that the asset allocation for a minimum variance portfolio
with expected return $\mu_P$ can also be written as follows in terms of the
asset weights vector $w_T$ of the tangency portfolio:
\begin{itemize}
\item asset weights vector $w_{min}$ given by
\begin{equation}
w_{min} \;=\; \frac{\mu_P - r_f}{\bar\mu^t w_T}\,w_T.
\end{equation}
\item weight $w_{min,cash}$ of the cash position given by
\begin{equation}
w_{min,cash} \;=\; 1 - \bfone^t w_{min};
\end{equation}
\end{itemize}
(ii) Write a pseudocode for computing the asset allocation for a minimum
variance portfolio with expected return $\mu_P$ using the formulas (9.102)
and (9.103).

\item (i) Show that the asset allocation for a maximum return portfolio with
variance of return equal to $\sigma_P^2$ can also be written as follows in
terms of the asset weights vector $w_T$ of the tangency portfolio:
\begin{itemize}
\item asset weights vector $w_{max}$ given by
\begin{equation}
w_{max} \;=\; \frac{\sigma_P}{\sqrt{w_T^t\Sigma_R w_T}}\cdot\operatorname{sign}\bigl(\bfone^t\Sigma_R^{-1}\bar\mu\bigr)\,w_T;
\end{equation}
\item weight $w_{max,cash}$ of the cash position equal to $1-\bfone^t
w_{max}$, i.e.,
\begin{equation}
w_{max,cash} \;=\; 1-\bfone^t w_{max}.
\end{equation}
\end{itemize}
(ii) Write a pseudocode for computing the asset allocation for maximum
return portfolio with variance of return equal to $\sigma_P^2$ using the
formulas (9.104) and (9.105).

\item Assume that the 99\% two day VaR of a portfolio is \$10 million.
Estimate the two day 95\% VaR of the portfolio and the five day 95\% VaR of
the portfolio.

\item Show that if two portfolios and their combined portfolio have normally
distributed returns, then the VaR of the combined portfolio is smaller than
the sum of the VaRs of each individual portfolio.

In other words, if $V_1$, $V_2$, and $V = V_1+V_2$ are portfolios with normal
returns, show that
\[
\VaR_V(N,C) \;\le\; \VaR_{V_1}(N,C) + \VaR_{V_2}(N,C),
\]
where $\VaR_V(N,C)$, $\VaR_{V_1}(N,C)$, and $\VaR_{V_2}(N,C)$ are the $N$ day
$C\%$ VaRs of $V$, $V_1$, and $V_2$, respectively.

Hint: Recall that, for a portfolio $W$ with normal returns,
\[
\VaR_W(N,C) \;=\; \sqrt{\frac{N}{252}}\,\sigma_{R_W} z_C W(0) \;-\; \frac{N}{252}\mu_{R_W}W(0),
\]
and use the fact that the return $R_V$ of the combined portfolio is the
weighted average of the returns $R_{V_1}$ and $R_{V_2}$ of the two
portfolios, i.e.,
\[
R_V \;=\; \frac{V_1(0)}{V(0)}R_{V_1} \;+\; \frac{V_2(0)}{V(0)}R_{V_2}.
\]

\item Consider three assets with multivariate normal distribution of their
returns. The expected values, standard deviations, and correlations of their
returns are
\[
\begin{aligned}
\mu_1 &= 0.08; & \sigma_1 &= 0.25; & \rho_{1,2} &= -0.25;\\
\mu_2 &= 0.12; & \sigma_2 &= 0.25; & \rho_{2,3} &= -0.25;\\
\mu_3 &= 0.16; & \sigma_3 &= 0.30; & \rho_{1,3} &= 0.25.
\end{aligned}
\]
\begin{subparts}
\item What is the 5--day 95\% VaR of \$100 million portfolios fully invested
in the first asset, fully invested in the second asset, and fully invested
in the third asset, respectively?
\item What is the minimal 5--day 95\% VaR of a \$100 million portfolio fully
invested in the first and second asset, fully invested in the second and
third asset, and fully invested in the first and third asset, respectively?
\item What is the minimal 5--day 95\% VaR of a \$100 million portfolio fully
invested in all three assets?
\end{subparts}

\end{exercises}
